% temporary newpage to make easier to write :)
\newpage


\stepcounter{section}
\section*{Lecture 2: details of sheaves and ringed spaces}

% TODO: 
% Remark about two ways to define manifolds, relate to schemes
% Define presheaf
% Examples presheaf
% Define sheaf
% Examples sheaf
% Remark sheaf with empty set
% Define ringed space
% Examples ringed spaces
% Define maximal ideal (???)
% Proposition about local ring / maximal ideal

\begin{remark}
    A manifold is a topological space equipped with extra data which forces it to ``locally resemble'' Euclidean space. Recall that there are two (equivalent) ways to define a manifold from a base topological space:
    \begin{enumerate}
        \item Using charts and atlases.
        \item By attaching a sheaf of functions.
    \end{enumerate}

    Usually one learns the chart definition and then derives and uses properties of sheaves while working with manifolds. We will work towards defining schemes in a similar manner, but primarily focus on the sheaf construction instead, and the analogues to charts will follow.
\end{remark}

We will need to define the machinery of sheaves to build towards schemes.


\begin{definition}[Presheaf]
    A \emph{presheaf} (of sets) $\mathcal{F}$ on a topological space $X$ consists of the following data:
    \begin{enumerate}
        \item For each open $U \subseteq X$, a set $\mathcal{F}(U)$.
        \item For each inclusion of open sets $V \subseteq U$, a \emph{restriction} function 
            \begin{align*}
            % Not sure what notation we're using so this for now , also should add map name to tikz diagram 
                \text{res}^U_V : \mathcal{F}(U) \to \mathcal{F}(V)
            \end{align*}
            which satisfies the following properties (contravariant functoriality):
            \begin{enumerate}
                \item For $V = U$, the restriction map for $V \subseteq U$ is the identity map $\text{res}_U^U = \text{id}_{\mathcal{F}(U)}$.
                \item For $W \subseteq V \subseteq U$, the following diagram commutes:
                    % https://q.uiver.app/#q=WzAsMyxbMCwwLCJcXG1hdGhjYWx7Rn0oVykiXSxbMiwwLCJcXG1hdGhjYWx7Rn0oVikiXSxbMiwyLCJcXG1hdGhjYWx7Rn0oVSkiXSxbMSwwXSxbMiwxXSxbMiwwXV0=
                    \[\begin{tikzcd}
                    	{\mathcal{F}(W)} && {\mathcal{F}(V)} \\
                    	\\
                    	&& {\mathcal{F}(U)}
                    	\arrow[from=1-3, to=1-1]
                    	\arrow[from=3-3, to=1-1]
                    	\arrow[from=3-3, to=1-3]
                    \end{tikzcd}\]
                
            \end{enumerate}
          \end{enumerate}
   A presheaf of groups/abelian groups/rings is a presheaf \(\mathcal{F}\) where \(\mathcal{F}(U)\) are groups/abelian groups/rings and where the restriction maps are homomorphisms.          
\end{definition}

We call the elements of \(\mathcal{F}(U)\) \emph{sections of \(\mathcal{F}\) on \(U\)} and the elements of \(\mathcal{F}(X)\) the \emph{global sections} of \(\mathcal{F}\).
For $f \in \mathcal{F}(U)$ and $V \subseteq U$, the restriction $\text{res}_V^U (f)$ is typically written  as $f|_V$, in line with restriction of functions.



\begin{example}\label{ex:presheaves}
    Let $X$ be a topological space. The following are examples of presheaves.

    \begin{enumerate}
    \item Take $X = \R^n$ with the usual Euclidean topology.
      Let \(\mathcal{F}\) be the presheaf whose sections on an open set \(U\) are the smooth functions \(f \colon U \to \mathbb{R}\), and the restriction maps are the usual restriction maps.
    \item 
       Let \(\mathcal{C}_X\) be the presheaf whose sections on an open set \(U\) are the continuous functions \(f \colon U \to \mathbb{R}\).
     \item This is a generalisation of the example above.  Replace the co-domain \(\mathbb{R}\) by any other topological space \(Y\).
     \item Suppose we have a topological space \(Y\) with a map \(\pi \colon Y \to X\).
       Let \(\operatorname{Sect}_{\pi}\) be the presheaf whose sections on \(U\) are continuous functions \(f \colon U \to Y\) such that \(\pi \circ f = \operatorname{id}_{U}\).
       This is called the sheaf of sections of \(\pi\).
              
        \item Fix a set $S$. For each open $U \subseteq X$, define $\mathcal{F}(U) = S$, and the restriction map to be $\text{id}_S$. This is the \emph{constant presheaf}.
        % \item Fix a topological space $Y$. The presheaf of maps into $Y$ is denoted
        %     \[
        %         \text{Maps}(-, Y) 
        %     \]
        % and let $U \subseteq X$ be open. Then the set of maps $f: U \to Y$, denoted $\text{Maps}(-,Y)$
        % \item 
    \end{enumerate}
    
\end{example}

% wanted to include footnote in case people forgot from analysis and for easy reference? probably not necessary.
\begin{definition}[Sheaf]

    Let \(X\) be a topological space and let  $\mathcal{F}$ be a presheaf on \(X\).
  We say that $\mathcal{F}$ is a \emph{sheaf} if for every open subset \(U \subset X\) and an open covering of \(U\) by open sets \(\{U_i\}_{i \in I}\), the following compatibility conditions are satisfied:
    \begin{enumerate}
        \item \textbf{Locality.} If $f, g \in \mathcal{F}(U)$ satisfy $f|_{U_i} = g|_{U_i}$ for all $i \in I$, then $f = g$. That is, ``sections are equal if they are locally equal''.
        \item \textbf{Gluing.} Given $f_i \in \mathcal{F}(U_i)$ such that $f_i|_{U_i \cap U_j} = f_j |_{U_i \cap U_j}$ for all $i, j \in I$, there exists a (unique) $f \in \mathcal{F}(U)$ such that $f|_{U_i} = f_i$ for all $i \in I$. 
    \end{enumerate}
\end{definition}

\begin{remark}
    The gluing condition only requires the existence of a section which restricts correctly to the open cover. However, uniqueness of this section is guaranteed by the locality condition.
\end{remark}

\begin{example} All of the presheaves in \Cref{ex:presheaves} are also examples of sheaves, except the last one (the constant presheaf).  If \(S\) contains more than one element, then the constant presheaf is not a sheaf (see the following proposition). 
\end{example}

However, not all presheaves are sheaves. The following is an important property of sheaves, which easily excludes certain presheaves from being sheaves.

\begin{proposition}[Sheaf behaviour on the empty set]
    If $\mathcal{F}$ is a sheaf, then $\mathcal{F}(\varnothing) = \{\cdot\}$ is a singleton set. 
\end{proposition}

\begin{proof}
    Let $f, g \in \mathcal{F}(\varnothing)$. The only open cover of $\varnothing$ is the empty covering. So $f = g$ (vacuously) on every part of the open cover. Then by the locality condition, $f = g$. So $\mathcal{F}(\varnothing)$ has exactly one element.
\end{proof}



\begin{nonexample}
    Fix $S$ to be a non-singleton set. The constant presheaf with $\mathcal{F}(U) = S$ is not a sheaf, because $\mathcal{F}(\varnothing) = S$ which is not a singleton.
\end{nonexample}

One might think we can fix this by manually setting $\mathcal{F}(\varnothing)$ to some singleton. But even this does not work in general, because it can break the gluing condition.

\begin{nonexample}
    Fix $X = \R$ and $S = \Z$. Take the constant presheaf, but modify it so that now $\mathcal{F}(U) = \{0\}$ (chosen arbitrarily to satisfy the singleton condition). In particular, take $U_1 = (0,1)$ and $U_2 = (2,3)$, which are disjoint and form an open cover for $U = U_1 \cup U_2$. Suppose $f_1 \in \mathcal{F}(U_1)$ is the constant function $1$, and $f_2 \in \mathcal{F}(U_2)$ is constant $-1$. These functions are compatible as they both restrict to the same element, $0$, on their intersection. But there is no function $f \in \mathcal{F}(U)$ which restricts correctly on both $U_1$ and $U_2$, because $1 \neq -1$.
\end{nonexample}

% It is also common for the topology of the underlying space to be too "non-algebraic" to admit 




% include something about locally constant being a valid sheaf?

Now we come to the defintion of a ringed space, which will be key moving forwards.

\begin{definition}[Ringed space]
    A \emph{ringed space} is a topological space $X$, equipped with a sheaf of rings $\mathcal{O}_X$.
\end{definition}

\begin{example}
  The topological space \(X = \mathbf{R}^n\) with \(\mathcal{O}_X = \mathcal{C}_X\) (the sheaf of continuous functions) is a ringed space.
  So is \(X = \mathbf{R}^n\) with \(\mathcal{O}_X = \mathcal{C}^{\infty}_X\) (the sheaf of smooth functions).
\end{example}

In future lectures we will be able to take a ring $R$, and construct from it a ringed space called $\text{spec}(R)$. Each $\text{spec}(R)$ will be an affine scheme, which we can glue together to form general schemes.



% An important property of ringed spaces is that they are entirely informed by the algebraic structure of a ring. 

\begin{nonexample}
    It is tempting to consider as an example, the space $X = \R$ with the presheaf as $\mathcal{F}(U) = \{\text{polynomial functions} : U \to \R \}$. But the gluing condition can break, for example if we take the zero polynomial on $U_1 = (-\infty, 0)$, and $x^2$ on $U_2 = (0, \infty)$. The polynomials automatically agree on the (empty) intersection, so should be compatible. But there is no global polynomial on the union $\R \setminus \{0\}$ which correctly restricts onto both $U_1$ and $U_2$.
\end{nonexample}

So what has gone wrong here, why is this not a sheaf? The answer is in the topology of the base space. The standard topology on $\R$ is the Euclidean topology, which relies on a number of non-algebraic operations such as square rooting, that may not be available or even compatible with rings. This motivates defining a topology which will be nicely compatible with all the algebraic properties we want to use. That will be another day. 

For now, let's build up some terminology to work with sheaves.

\begin{definition}[Stalk]
    Let $(X, \mathcal{F})$ be a sheaf. For a given point $x \in X$, we define the \emph{stalk} of $\mathcal{F}$ at $x$ to be 
    \begin{align*}
        \mathcal{F}_x & = \varinjlim_{U \ni x} \mathcal{F}(U) \\
        & = \bigsqcup_{U \ni x, f \in \mathcal{F}(U)} (U,f) \slash \sim
    \end{align*}
    where $(U,f) \sim (V, g)$ if there is a $W \ni x$, $W \subseteq U \cap V$, such that $f|_W = g|_W$. These equivalence classes are called \emph{germs}.
  \end{definition}

  If \(\mathcal{F}\) is a sheaf of groups, then the stalk of \(\mathcal{F}\) at a point \(x\) is a group.
  Similarly, for abelian groups, rings, etc.

The stalk captures the behaviour of the sheaf around a given point. 
\begin{example}
  Let $X = \R$ and $\mathcal{F} = \mathcal{C}_X$ the sheaf of continuous functions on $X$.
  The stalk $\mathcal{F}_0$ is a ring.
  It contains a copy of $\R$, corresponding to the constant functions, but it also contains other elements, represented by functions like $x, x^2, x^3, \sqrt{x}, e^x, \sin(x), |x|$ (all these represent distinct germs).

  The ring \(\mathcal{F}_0\) has no nilpotents, but is not an integral domain.
  It is a ``local ring'', meaning it has exactly one maximal ideal.
  The maximal ideal is given by the kernel of the evaluation map at zero:
    \[
        m = \ker (\text{eval}_0) = \{f \; | \; f(0) = 0\}.
    \]
    The evaluation map at zero is surjective onto $\R$ (consider the constant functions).
    So, the first isomorphism theorem gives $\mathcal{F}_0 \slash m \cong \R$.
    Since $\R$ is a field, this means $m$ is maximal.
    To see that \(m\) is the unique maximal ideal, we see that every element of \(\mathcal{F}_0\) that is not contained in \(m\) is a unit.
    Indeed, if \(f \in \mathcal{F}_0\) is such that \(f(0) \neq 0\), then \(f(x) \neq 0 \) for all \(x\) in a neighbourhood \(U\) of \(0\).
    Then the germ \((U, 1/f)\) is the multiplicative inverse of \(f\).
    Since the complement of \(m\) in \(\mathcal{F}_0\) consists of units, every non-unit ideal of \(\mathcal{F}_0\) must be contained in \(m\).
    So \(m\) is the only maximal ideal.
  \end{example}
  

% \begin{remark}
%     Arbitrary topologies (such as the standard Euclidean topology) may not behave according to the rigid properties of algebraic structures such as rings. 
% \end{remark}




% For now, let's become familiar with ringed spaces.

%%% Local Variables:
%%% mode: LaTeX
%%% TeX-master: "main"
%%% End:
